\section*{Hoofdstuk \ref{ch:b2:plant-meristems} --- Vegetatieve ontwikkeling van planten: meristemen en groei}

\begin{solution}{exo:b2:plant-meristems:1}
Een \omterm{def:b2:plant-meristems:meristem}{meristeem} is een blijvende populatie ongedifferentieerde, delende
cellen waaruit de weefsels van de plant voortkomen. \omterm{def:b2:plant-meristems:meristem}{Apicaal
scheutmeristeem}: bladeren, stengel, \omterm{def:b2:angiosperm-reproduction:flower}{bloemen}; \omterm{def:b2:plant-meristems:meristem}{apicaal wortelmeristeem}: de
wortel; \omterm{def:b2:plant-meristems:wood}{vaatcambium}: secundair xyleem (\omterm{def:b2:plant-meristems:wood}{hout}) en floëem; kurkcambium: het
\omterm{def:b2:plant-meristems:wood}{periderm} (kurk) van de schors.
\end{solution}

\begin{solution}{exo:b2:plant-meristems:2}
Wortelmutsje (bescherming, smering, waarneming van de zwaartekracht,
afgeschuurde cellen); \omterm{def:b2:plant-meristems:meristem}{meristeem} met het rustende centrum (deling);
strekkingszone (cellen worden tien- tot twintigmaal langer);
differentiatiezone (wortelharen, rijp xyleem en floëem). Zijwortels
ontstaan uit de pericykel van de rijpe wortel, tegenover de xyleempolen.
\end{solution}

\begin{solution}{exo:b2:plant-meristems:3}
Bij bonenplanten werd de top weggesneden: de zijknoppen liepen uit. Bij een
tweede reeks kreeg de afgesneden stomp een agarblokje met auxine: de
knoppen bleven in rust; bij de controlereeks een blokje gewone agar: de
knoppen liepen uit. Conclusie: auxine uit de top is wat de zijknoppen remt
--- \omterm{prop:b2:plant-meristems:apical}{apicale dominantie} is hormonaal.
\end{solution}

\begin{solution}{exo:b2:plant-meristems:4}
Zestig jaar oud; het spinthout bestaat uit de laatste vijftien ringen;
levend zijn het \omterm{def:b2:plant-meristems:meristem}{cambium}, het floëem, de straal- en parenchymcellen van het
spinthout, en het kurkcambium --- een dunne schil rond een dode kern.
\end{solution}

\begin{solution}{exo:b2:plant-meristems:5}
$0.1\times(P - 0.3)$: \qty{0.01}{h^{-1}} (verdubbeling in $\ln 2/0.01 =
\qty{69}{h}$), \qty{0.03}{h^{-1}} (\qty{23}{h}), \qty{0.05}{h^{-1}}
(\qty{14}{h}). Bij \qty{0.25}{MPa} ligt de turgor onder de zwichtdrempel:
de groei stopt, nog vóór elke verwelking.
\end{solution}

\begin{solution}{exo:b2:plant-meristems:6}
Blad 1--8 op 0, 137.5, 275, 52.5, 190, 327.5, 105, \qty{242.5}{\degree}.
Gesorteerd: 0, 52.5, 105, 137.5, 190, 242.5, 275, 327.5 --- de kleinste
opening is \qty{32.5}{\degree}. Elk blad zit in een opening tussen de
vorige, zodat de acht bladeren elkaar zo weinig mogelijk beschaduwen en de
kroon de lichtschijf gelijkmatig vult.
\end{solution}

\begin{solution}{exo:b2:plant-meristems:7}
$2\pi r h\Delta r$: bij \qty{5}{cm} $2\pi\times 0.05\times 12\times
0.003 = \qty{0.011}{m^3}$, \qty{5.7}{kg}, \qty{2.8}{kg} koolstof; bij
\qty{40}{cm} \qty{0.090}{m^3}, \qty{45}{kg}, \qty{23}{kg} koolstof
--- acht keer meer bij dezelfde ringbreedte.
\end{solution}

\begin{solution}{exo:b2:plant-meristems:8}
\emph{clv3}: geen rem --- enorme meristemen, extra bladeren en bloemorganen,
gefascieerde stengels. \emph{wus}: geen gaspedaal --- het \omterm{def:b2:plant-meristems:meristem}{meristeem} raakt na
enkele bladeren uitgeput en wordt steeds opnieuw gesticht. Dubbelmutant: het
\emph{wus}-fenotype, dus werkt WUS stroomafwaarts van CLV3: de taak van CLV3
is WUS te onderdrukken, en zonder WUS valt er niets te onderdrukken.
\end{solution}

\begin{solution}{exo:b2:plant-meristems:9}
Transport: \qty{20}{h}. Diffusie: $(0.2)^{2}/(2\times 10^{-9}) =
\qty{2e7}{s}$, ongeveer 230 dagen. Polair transport maakt een hormonaal
signaal binnen een dag bruikbaar over de lengte van een plant, en geeft het
een richting.
\end{solution}

\begin{solution}{exo:b2:plant-meristems:10}
\qty{3}{cm} rijp weefsel per dag met cellen van \qty{200}{\micro m}:
\num{150} cellen per rij per dag. Honderd delende cellen moeten zich dus elk
1.5 keer per dag delen: een cyclus van ongeveer \qty{16}{h}.
\end{solution}

\begin{solution}{exo:b2:plant-meristems:11}
Gibberelline verlengt internodia; een plant die er niet op kan reageren,
heeft korte internodia en een korte stengel. Zware stikstofbemesting maakt
een hoge tarwe hoger en zwaarder van aar tot ze omvalt (legert) en rot; een
korte stevige stengel blijft staan en steekt de extra koolstof in graan
in plaats van stro, zodat de oogstindex stijgt. In het wild wordt een korte
plant door haar buren overgroeid en beschaduwd, en verliest ze.
\end{solution}

\begin{solution}{exo:b2:plant-meristems:12}
Meristemen laten een plant organen toevoegen waar en wanneer de
omstandigheden gunstig zijn --- bladeren naar het licht, wortels naar het
water --- en organen opgeven die niet renderen; elke okselknop is een
reservegroeipunt, zodat begrazing, vuur en snoei door hergroei worden
overleefd; en omdat de meristemen voortdurend embryonaal zijn, heeft een
\omterm{def:b2:asexual-reproduction:clone}{genet} geen vaste levensduur. Het nadeel: een plant kan wat ze heeft gebouwd
niet verplaatsen of reorganiseren --- een eenmaal geplaatst orgaan blijft
waar het is, en een lichaam dat door aangroei is gebouwd, moet zijn dode
verleden (het kernhout) meedragen.
\end{solution}

\begin{solution}{pb:b2:plant-meristems:1}
\textbf{1.} $150/3 = 50$ bladeren per top; \num{10000} voor de boom.
\textbf{2.} 0, 137.5, 275, 52.5 en \qty{190}{\degree}.
\textbf{3.} $8\times 137.5 = 1100 = 3\times 360 + 20$; $13\times 137.5
= 1787.5 = 5\times 360 - 12.5$; $21\times 137.5 = 2887.5 = 8\times 360
+ 7.5$: het 22e blad is het eerste binnen \qty{10}{\degree} van het
eerste.
\textbf{4.} Elk nieuw blad zit in de breedste opening, zodat de bladeren
elkaar het minst beschaduwen; het is auxine, waarvan het polaire transport
(PIN-dragers) de omgeving van elk primordium leegzuigt, dat de afstand
bepaalt.
\textbf{5.} $\num{10000}\times 80\,\text{cm}^{2} = \qty{80}{m^2}$: de
hele kroon bestaat uit bladeren van dit seizoen.
\textbf{6.} Alles; de eerste bladeren worden gebouwd uit zetmeel dat het
vorige jaar in het \omterm{def:b2:plant-meristems:wood}{hout} en de wortels werd opgeslagen.
\textbf{7.} Een verdrievoudigde top maakt sneller primordia --- twee of
drie per plastochron --- zodat bladeren en bladoppervlak zouden kunnen
verdrievoudigen, maar de scheuten zouden gefascieerd zijn, met lintvormige
stengels en een verstoorde fyllotaxis, en een groot deel van het extra blad
zou zichzelf beschaduwen.
\textbf{8.} Overdag: $0.05\times(0.55 - 0.35) = \qty{0.01}{h^{-1}}$;
's nachts: $0.05\times 0.40 = \qty{0.02}{h^{-1}}$.
\textbf{9.} Overdag: $2e^{0.12} = \qty{2.25}{cm}$, \qty{0.25}{cm}
toegevoegd; 's nachts: $2.25e^{0.24} = \qty{2.87}{cm}$, \qty{0.61}{cm}
toegevoegd.
\textbf{10.} Gemiddelde snelheid \qty{0.015}{h^{-1}} over \qty{96}{h}: $2e^{1.44}
= \qty{8.4}{cm}$.
\textbf{11.} Overdag: $0.05\times 0.30 = \qty{0.015}{h^{-1}}$; 's nachts:
$0.05\times 0.50 = \qty{0.025}{h^{-1}}$.
\textbf{12.} $0.05\times 0.05 = \qty{0.0025}{h^{-1}}$, een kwart van de
normale snelheid overdag. Osmotische aanpassing verhoogt het gehalte aan
opgeloste stoffen van de cel, zodat de turgor bij dezelfde waterpotentiaal
weer boven de drempel uitkomt.
\textbf{13.} Overdag verlaagt de transpiratie de waterpotentiaal van het
blad en daarmee de turgor, zodat $P - Y$ klein is; 's nachts wordt de
turgor hersteld. De groei wordt uur na uur beperkt door water, niet door
suiker.
\textbf{14.} $2\pi\times 0.10\times 8\times 0.005 = \qty{0.025}{m^3}$;
\qty{12.6}{kg} droog.
\textbf{15.} \qty{6.3}{kg} koolstof.
\textbf{16.} $80\times 4\times 150 = \qty{48}{kg}$ koolstof.
\textbf{17.} $6.3/48 = \qty{13}{\%}$. Andere bestemmingen: de bladeren
zelf, wortels, takken en twijgen, de ademhaling van elke levende cel
(ongeveer de helft van wat wordt vastgelegd), reserves, \omterm{def:b2:angiosperm-reproduction:flower}{bloemen} en zaden.
\textbf{18.} Straal \qty{0.20}{m}: $2\pi\times 0.20\times 8\times
0.005 = \qty{0.050}{m^3}$, twee keer die van dit jaar --- de omtrek
waarlangs het \omterm{def:b2:plant-meristems:meristem}{cambium} werkt, is verdubbeld (en de boom zal hoger zijn).
\textbf{19.} Een droog of koud jaar; of ontbladering door insecten, vorst of
vuur. Een klimatologische oorzaak versmalt dezelfde ringen in elke boom van
de streek en correleert met weergegevens; schade treft één boom of één
bestand, en vorst of vuur laat anatomische littekens in de ring achter.
\textbf{20.} De knoppen die het dichtst bij elke snede zitten, lopen binnen
enkele dagen uit tot meerdere scheuten; na een maand is de boom
struikvormiger en dichter, met meer maar kortere scheuten.
\textbf{21.} Auxine op de snede vervangt het signaal van de top: de knoppen
blijven in rust en er volgt geen vertakking.
\textbf{22.} Elke knipbeurt verwijdert de toppen en laat de knoppen eronder
los; zes keer per jaar voegt elke loslating een laag korte scheuten toe, en
het oppervlak raakt ermee gevuld.
\textbf{23.} Uit slapende okselknoppen op de stobbe en uit adventieve
knoppen die door het \omterm{def:b2:plant-meristems:meristem}{cambium} worden gevormd; elke knop is een volledig
\omterm{def:b2:plant-meristems:meristem}{meristeem} dat het scheutstelsel kan herbouwen, terwijl de kop van een dier
één enkele, onvervangbare \omterm{def:b2:vertebrate-development:induction}{organisator} is zonder reservekopieën.
\textbf{24.} Meer cytokinine dat de knoppen bereikt, bevordert hun
uitlopen --- een gesnoeide plant die water krijgt, vertakt zich meer.
\textbf{25.} 50 bladeren per top, \num{10000} per boom; groeisnelheden van
het internodium \qty{0.01}{h^{-1}} overdag en \qty{0.02}{h^{-1}} 's nachts;
het \omterm{def:b2:plant-meristems:wood}{hout} van het jaar \qty{0.025}{m^3}, \qty{12.6}{kg} droog, \qty{6.3}{kg}
koolstof.
\end{solution}
